% ===== §2 The Field of Travel and Its Constitutive Features =====
\section{The Field of Travel and Its Constitutive Features}
\label{sec:field}

\noindent
If travel creates a special field, the field can be analysed. This section does so, decomposing the field into the features that constitute it and arguing, for each, both what it is phenomenologically and how it bears on the relation between two people. The features fall into three groups. The first concerns the field's \emph{temporal arch}, the structure of the journey in time, which begins before departure and continues after return. The second concerns the field's \emph{openness}, what lifts the everyday's closure and makes the field a field at all. The third concerns the \emph{conditions of being-together}, what makes the field's openness specifically shared in place of the openness two solitary travellers might each enjoy. Underlying all three groups is a condition that the section treats last and separately, since it is the ground of their possibility in place of one feature among the others: spaciousness. The section closes by examining how the features interact, reinforce, and conflict, so that the field is understood as a dynamic system whose internal tensions are themselves a source of what travel does to a relation, in place of a static list of properties.

% ============================================================
\subsection{The temporal arch I: co-creation}
\label{sub:field-cocreation}

A journey does not begin at the airport. It begins when two people first imagine it together, when the idea is floated, the map opened, the possibilities argued over, the expectations aligned, the shared anticipation built. This phase of \emph{co-creation} is already the field itself, and indeed in one respect its purest stretch, in place of the preparatory labour that precedes the journey proper, on which the consumerist understanding takes planning to be merely the front-loaded work of arranging a product to be later consumed. In co-creation the two people jointly build a shared world that does not yet exist, and the building of a shared world is the foundational act by which a ``we'' is constituted.

The point is worth dwelling on, because it locates the generativity of travel earlier than the usual account places it. To plan a journey together is to negotiate a \emph{shared intention}: where to go, why, what to seek, what to refuse, how to weigh the one's wish to see against the other's wish to rest. A shared intention is a single intentionality borne by two rather than the sum of two private intentions, and its construction is the forging of the relational subject in one of its most legible forms. The everyday offers surprisingly few occasions for this kind of pure joint construction, since routine consists largely in the discharge of intentions already given, the roles, the duties, the standing arrangements, in place of the building of new ones together. Travel planning forces it. Two people imagining a place they have not yet been, aligning what they each hope to find there, are doing in miniature exactly what it is to be a ``we'': they are jointly intending a world. And because what they jointly intend is a future background, the manifold against which their resonance will, in a few weeks, be redefined, co-creation is, quite precisely, the anticipatory and conspiratorial rehearsal of the background-exchange that the journey will enact. Before the manifold is exchanged in fact, the two have already begun, in imagination, to build it together. Co-creation is the first opening of the field, and because it recurs, in the re-planning of each day en route and, as the paper will argue, in the shared retelling afterward (\xref{sec:reproduction}), it forms the axis along which the whole temporal arch is strung more than a single phase.

% ============================================================
\subsection{The temporal arch II: reshaped time}
\label{sub:field-time}

Travel reshapes time. The everyday's time is, in the strict sense, \emph{told off} into units assigned to functions: the hour for the commute, the hour for the meeting, the evening portioned among its obligations, a temporality whose every interval is already spoken for. The time of a journey has a different texture, and the difference exceeds its being leisure in place of labour. It is, first, \emph{framed}: a journey has a beginning and an end, a definite span lifted out of the indefinite flow of ordinary time and bounded as a whole. This framing gives the time of travel something of the quality of an exception, a parenthesis in the run of the year, and the exceptional, bounded character contributes to its value. Precisely because the journey will end, because its time is finite and counted, each moment within it is held more closely; the knowledge that one is in a span that will not last presses one into a fuller presence within it than the seemingly endless, because unframed, time of the everyday ever compels.

It is, second, \emph{continuous and shared}. The everyday distributes two people's time into largely separate channels that intersect at the day's margins, the breakfast, the evening, the weekend's negotiated overlap, so that much of a shared life is in fact lived in parallel more than together. The time of a journey is unbroken joint time: hour upon hour side by side, with no separate offices to disperse into, no independent circles to absorb the day. This continuity is one of the conditions that makes the field's other effects possible, since many of them, the forcing of co-presence, the reaching of depth, require a sustained, uninterrupted being-together that the channelled time of the everyday structurally prevents. Reshaped time, framed and continuous, is the temporal medium within which the rest of the field can do its work; and together with co-creation it completes the field's temporal structure, the arch that runs from the first imagining, through the bounded continuous span of the journey itself, to the retelling that will reopen the span in memory.

% ============================================================
\subsection{The openness I: suspension}
\label{sub:field-suspension}

The everyday's closure is lifted, first, by \emph{suspension}: the bracketing, for the duration of the journey, of the dense network of roles and obligations that ordinarily defines each person. At home one is the one who must go to work, the one who owes the mortgage, the one who must manage the relatives, the one assigned this place in this set of standing arrangements; and these assignments are, in large part, constitutive of how one shows up day to day, more than external decorations upon a self that would be the same without them. To travel is to bracket them, suspending their operation in place of abolishing them, since they wait at home, so that for a span one is, more nearly, simply oneself and simply with the other in place of the discharger of those functions. The structure is close to what phenomenology calls the \emph{epoché}: the methodical bracketing of the natural attitude so that what was hidden by its automatic operation can be seen. Here the bracketing is performed by a plane ticket in place of a philosopher's resolve, and what it makes visible is the two people underneath their roles, who they are, and who they are to each other, when the scaffolding of function is set aside.

This is among the reasons travel is at once a touchstone and a danger. The roles that suspension brackets are, for many couples, a good part of what has been holding the relation in working order; bracket them, and the relation must stand, for a while, on whatever lies beneath. If beneath the roles there is a living mutuality, suspension liberates it, and the couple discover a freedom in being together that the role-bound everyday had cramped. If beneath the roles there is little, suspension exposes the little, and the couple discover, sometimes with a vertigo they would rather not have had, how much of their togetherness was the interlocking of functions and how little was anything else. Suspension removes what was concealing what it reveals in place of creating it. That is the precise sense in which the field is a touchstone: it assays the relation by taking away the supports and seeing what stands.

% ============================================================
\subsection{The openness II: contingency}
\label{sub:field-contingency}

The second source of the field's openness is \emph{contingency}: the lifting of the everyday's predetermination, the restoration to each moment of the quality of being genuinely undecided. The everyday is, in its essence, repetition, and repetition is the predetermined: one knows, in the ordinary run of a day, more or less what each hour will hold, because each hour has been held a thousand times before. This predetermination is what the existential tradition has taught us to recognise as the great anaesthetic of selfhood. In Heidegger's analysis, the everyday is the reign of \emph{das Man}, the anonymous ``they'' in whose averaged, levelled-down understanding existence is absorbed: one does what one does, as ``one'' does it, and in this absorption the fact of one's own existence as a question to be lived is covered over, fallen into the smooth automatism of the familiar \citep{heidegger1962b}. Travel breaks the repetition, and the breaking is principally a restoration of the undecided more than a matter of fresh stimulation. In an unfamiliar place nothing is settled in advance: one does not know whether the train will come, where the lane leads, whether one will be understood, whether the food will be good. Each moment becomes, again, a genuine present, undetermined, open, to be undergone in place of discharged.

And in this restored contingency the subject is pressed back into appearing. When the familiar scaffolding is removed, the fact that \emph{I am here, undergoing this, having to choose} becomes visible again, no longer dissolved into the automatism of the predetermined. This is the existential core of travel, and it has a long and serious philosophical lineage, which the field's contingency reactivates in place of inventing. The lineage is summarised in Table~\ref{tab:existential}. The point common to its members, however they differ, is that the rupture of the predetermined returns the subject to itself, to its existence, its freedom, its finitude, by stripping away the everyday's concealment. What this paper adds to that lineage is the move from the solitary to the shared, taken up in the next section; here it is enough to establish that contingency is one of the two great openers of the field, and that its opening is, in the first instance, existential, the return of existence to presence.

\begin{table}[ht]
\centering
\small
\begin{tabularx}{\linewidth}{L{3.0cm} L{3.3cm} Y}
\rowcolor{tablehead}\lrhead{Thinker} & \lrhead{Concept} & \lrhead{Bearing on the contingency of the field} \\
\midrule
Heidegger \citep{heidegger1962b} & \emph{das Man}, fallenness, authenticity; anxiety (\emph{Angst}) as disclosive & The everyday's repetition is absorption in the ``they''; contingency lifts the concealment, a bearable anxiety in which existence is disclosed. The paper's principal mast. \\
Sartre \citep{sartre1964,sartre1956} & Contingency; the nausea before the gratuitousness of being; freedom & Travel restores the felt gratuitousness of things, their lack of necessary reason, and with it the exposure of freedom: in the undetermined moment, one must choose. \\
Kierkegaard \citep{kierkegaard1983} & Repetition vs.\ recollection; genuine ``repetition'' as renewed receiving & Sharpens the paper's own wording: the aim is the renewed winning of existence out of contingency rather than mechanical re-enactment. \\
Jaspers \citep{jaspers1970} & Boundary situations (\emph{Grenzsituationen}) & The journey's mishaps, getting lost, stranded, stuck, are minor boundary situations in which ordinary self-understanding fails and the subject meets its finitude and possibility. \\
Merleau-Ponty \citep{merleauponty2012} & The lived body; perception as bodily engagement with a world & The unfamiliar place is met first in the body, disoriented, alert, re-learning its surroundings, so that contingency is undergone perceptually before it is understood. \\
\end{tabularx}
\caption{The existential-phenomenological lineage reactivated by the field's contingency. The paper takes Heidegger as its principal mast and reads the others as developing aspects of the single thesis that the rupture of the predetermined returns the subject to presence.}
\label{tab:existential}
\end{table}

% ============================================================
\subsection{The conditions of being-together I: liminality}
\label{sub:field-liminality}

The features so far would each be available to a solitary traveller. The next three are what make the field's openness specifically \emph{shared}. The first is \emph{liminality}. The concept is anthropological: in van Gennep's analysis of rites of passage, and Turner's development of it, the rite has a threshold phase, the \emph{limen}, in which the initiate has left the old status without yet entering the new, and is for a time betwixt and between, outside the ordinary structure of roles and ranks \citep{vangennep1960,turner1969}. Two things characterise this liminal phase. The first is the temporary dissolution of structure: in the threshold, the ordinary hierarchies and classifications are suspended, and the initiates stand outside the grid that ordinarily places them. The second is what Turner called \emph{communitas}: the intense, levelling, immediate bond that arises among those who pass through the liminal together, a being-with stripped of the mediations of status, an undifferentiated mutuality possible only in the threshold \citep{turner1969}.

Travel is a liminal condition, and shared travel generates a communitas of two. To be on a journey is to be betwixt and between, to have left the structure of home without having arrived at any new structure, suspended in a passage between two settled states. And two people who enter this threshold together are, for its duration, released from the structure that ordinarily mediates them, the division of household labour, the standing roles, the accreted positions of a shared life, into a more immediate and levelling mutuality. The couple on the road are not, for a while, the one-who-manages-the-finances and the one-who-handles-the-relatives; they are two people in the threshold together, met more nearly as such. This dyadic communitas is one of the deepest goods of shared travel, and it is also why return can be so deflating: communitas is, by its nature, a threshold phenomenon, and the re-entry into structure that homecoming requires (\xref{sec:reproduction}, \xref{sub:reproduction-return}) dissolves it unless deliberate labour preserves what it generated.

% ============================================================
\subsection{The conditions of being-together II: symmetric unfamiliarity}
\label{sub:field-unfamiliarity}

The second condition of shared openness is \emph{symmetric unfamiliarity}, and though it is easily overlooked it may be the most important of the three, because it is the condition under which genuine \emph{shared witnessing} becomes possible. In the everyday, familiarity is almost never symmetric. The home is the home of someone first; the city is one partner's native ground and the other's adopted one; the circle of friends originates with one and is acquired by the other; the domains of competence are divided, so that in each there is a host and a guest, one who knows and one who is shown. This asymmetry is not malign, but it is pervasive, and it means that in the ordinary run of a shared life the two are rarely, if ever, set before the same unknown on equal terms. Travel to a place strange to both removes the asymmetry. The foreign city is equally foreign to the two of them; neither is the host, neither holds the home advantage, neither knows the way. For once they stand before the same unknown as equals.

The significance of this is that it is the condition of \emph{co-witnessing}, as against the asymmetric witnessing of host and guest. When one partner shows the other their home city, the other witnesses the city, and also witnesses the partner-as-host, and is witnessed as guest; the structure has a giver and a receiver of the unfamiliar. When both are equally strangers, neither hosts and neither is hosted; they undergo the unknown \emph{together}, side by side, each seeing it freshly and each seeing the other see it freshly, and, the decisive addition, each knowing that the other is in the same condition. This symmetric co-undergoing is the ground of the field's deepest generativity, because it is the condition under which the relation's dynamics can become present to both at once as something they are equally inside of, in place of something one is showing and the other receiving. The next section's central concept, the self-presence of relational dynamics, depends upon this symmetry: it is hard to feel the ``we'' as a whole when one of the two is the host of the occasion; it becomes possible when both are equally in the threshold, equally before the same strangeness, equally without a script.

% ============================================================
\subsection{The conditions of being-together III: forced co-presence}
\label{sub:field-copresence}

The third condition is the most prosaic and the most double-edged: \emph{forced co-presence}. Travel puts two people together continuously, with the everyday's escape valves removed. At home there are always exits, the separate job, the independent friendships, the room one can withdraw to, the errand that takes one out, and these exits, though they look like mere logistics, perform a real function: they release the pressure that continuous proximity generates, and they let a relation run without ever having to test how much undiluted togetherness it can bear. Travel removes the exits. For the span of the journey the two are thrown together without reprieve, and this enforced, unrelieved proximity is at once the field's chief stressor and one of its chief engines of depth.

As stressor, it amplifies. The small frictions that the everyday's exits would have dispersed, the differing rhythms, the divergent wants, the accumulated irritations, have, on the road, nowhere to go; they cannot be walked off into the separate channels of a normal day, and so they concentrate, which forms a second part of the explanation, beyond suspension and contingency, of why couples quarrel on holiday. Yet the same removal of exits that amplifies friction also forecloses the easy avoidances by which a relation evades its own depths, and so forces a genuine co-presence that the exit-rich everyday permits one perpetually to defer. With nowhere to escape to, the two must actually be with each other, meeting whatever is between them in place of managing it, and in that enforced meeting a depth becomes reachable that the avoidant comforts of home keep forever at arm's length. Forced co-presence is, like suspension, a feature whose value is inseparable from its danger: it is the same pressure that cracks a brittle relation and tempers a sound one.

% ============================================================
\subsection{The underlying condition: spaciousness}
\label{sub:field-spaciousness}

The seven features so far do not, however, sit as seven coordinate items on a list. One condition underlies them all, and it is the ground that makes the rest possible in place of one feature among them: \emph{spaciousness}, the worked emptiness introduced in \xref{sub:intro-spaciousness}. Each of the seven, examined closely, turns out to require an opened, unfilled room in order to operate. Contingency cannot enter a filled field: a journey whose every hour is booked, whose every space is crowded with planned content, leaves no gap through which the undetermined can arrive, and the wildflower at the roadside cannot be seen when the eyes are fixed on the itinerary. Suspension is itself a kind of emptying, since the bracketing of roles is the clearing-away of the functions that fill the everyday. The liminal threshold is structurally an emptiness, the cleared interval between two structures. Forced co-presence does its generative work only where there is unfilled time for the two to actually occupy together; pack the journey densely enough and the co-presence becomes mere joint logistics, two people processing an itinerary side by side. Even co-creation and reshaped time are, at bottom, ways of opening and bounding a space that is to be left partly empty, a vessel whose use is in the room within.

This is why spaciousness is named as the field's underlying condition in place of its eighth feature. It is the emptiness in the vessel; the others are the walls that shape and hold it. The eleventh chapter of the \emph{Daodejing} states the structure exactly: what is there, the walls, the features, the bounded span, gives advantage, while what is not there, the emptiness they enclose, gives use \citep{laozi1963}. The advantage of the field lies in its constructed features; its \emph{use}, the actual generation of presence and of the ``we,'' lies in the emptiness those features enclose and protect. And this is the deepest reason the power of travel is misunderstood as novelty. Novelty fills; spaciousness empties; and the emptying, in place of the filling, opens the room in which a relation can come to itself. The phrase the previous paper used for the highest form of this, the fullness of doing nothing, two people before a window with the rain \citep{paperxix}, was already a description of spaciousness, here given its name and its place as the condition of the whole field. To travel well is, before all the particular features, to know how to keep the vessel empty enough to be of use.

% ============================================================
\subsection{The features in tension: the field as a dynamic system}
\label{sub:field-tensions}

It would falsify the field to leave its features as a static inventory, for they do not sit side by side in placid coexistence; they pull on one another, reinforce one another, and, crucially, conflict, and the field's actual character in any given journey is the resultant of these tensions rather than the sum of the features taken singly. Three of the tensions are worth drawing out, both because they are phenomenologically real and because they will matter to the ethics and the practice that follow.

The first is the tension between \emph{co-creation and contingency}, which forms the deep structure of a familiar quarrel. Co-creation, carried to excess, is the enemy of contingency: the more completely a journey is planned, every hour assigned, every site booked, every meal chosen in advance, the less room remains for the undetermined to enter, until the journey becomes a fixed background in motion, the everyday's predetermination carried abroad under the costume of travel. Yet contingency without any co-creation is a drift, and an anxious one, in place of a journey; some shared intending is the condition of there being a ``we'' that travels at all. The over-planner and the under-planner, the one who would book the journey to the minute and the one who would simply set out, are pulling on two genuinely constitutive features of the field rather than merely standing temperamentally at odds, each feature of which, alone and to excess, would destroy it. The resolution this paper will eventually offer (\xref{sec:construction}) is that the right object of planning is the \emph{frame} in place of the \emph{content}: that co-creation should build and protect a bounded, emptied space and then leave the interior open to contingency, intention exercised on the vessel's walls, the emptiness within left for what will come.

The second is the tension between \emph{forced co-presence and spaciousness}. Forced co-presence supplies the continuous joint occupation the field requires; but spaciousness requires that the joint time remain not wholly filled, including not wholly filled \emph{with each other}. A togetherness with no interior emptiness, no room for either to be quietly alone within the shared field, no unfilled silence, becomes a suffocation, and the relentless proximity that tempers a sound relation can, without spaciousness, smother it. The field needs both the removal of the exits and an internal room; co-presence without spaciousness is confinement, and spaciousness without co-presence is two people merely near one another.

The third is the tension between \emph{liminality and the temporal arch}, which is the structural seed of the problem of return. The communitas of the threshold is, by its nature, anti-structural and impermanent; the temporal arch, by its nature, ends, depositing the two back into the structure of the everyday. The very feature that makes the field's mutuality so intense, its liminal release from structure, is the feature that guarantees its dissolution upon return, since one cannot remain forever in the threshold. This tension is the condition that makes the labour of reproduction necessary, in place of a defect to be removed (\xref{sec:reproduction}): because the liminal cannot be prolonged, what it generated must be \emph{reproduced} into a form that can survive the return to structure, or it is lost. The field, then, is a dynamic and self-undermining condition rather than a stable possession, intense and impermanent, whose goods must be actively carried across the threshold of return or surrendered at it. It is to that carrying, and to the wealth it conserves or squanders, that the paper now turns, after first attending to what the field, while it lasts, brings to presence.
